\documentclass[tikz,border=4pt]{standalone}
\usepackage{ctex}
\usepackage{tikz}
\usetikzlibrary{shapes.geometric, arrows.meta, positioning}
\begin{document}
\begin{tikzpicture}[
    node distance=0.8cm and 1.4cm,
    box/.style={rectangle, rounded corners=3pt, draw=blue!70, fill=blue!10,
                text width=3.8cm, align=center, font=\small, minimum height=0.8cm},
    dec/.style={diamond, draw=orange!80, fill=orange!15, aspect=2.2,
                text width=3cm, align=center, font=\small},
    arr/.style={-Stealth, thick},
    yes/.style={font=\small, sloped, above},
    no/.style={font=\small, sloped, below},
]
\node[box] (start) {给定级数 $\sum a_n$};
\node[dec, below=of start] (nec) {$a_n\to 0$？};
\node[box, right=2.2cm of nec] (div0) {发散};
\node[dec, below=of nec] (pos) {正项级数？};
\node[dec, below left=1.2cm and 0.2cm of pos] (ratio) {$\rho=\lim\frac{a_{n+1}}{a_n}$？};
\node[dec, below right=1.2cm and 0.2cm of pos] (alt) {交错级数？};
\node[box, below=of ratio] (ratioR) {$\rho<1$收敛\\$\rho>1$发散\\$\rho=1$失效};
\node[box, below=of alt] (leibniz) {Leibniz:\\$|a_n|\downarrow 0$收敛};
\node[box, below=2cm of nec] (compare) {};
\node[box, below=1.2cm of ratioR] (final) {绝对收敛\\$\Rightarrow$ 收敛};

\draw[arr] (start) -- (nec);
\draw[arr] (nec) -- node[right, font=\small]{否} (div0);
\draw[arr] (nec) -- node[right, font=\small]{是} (pos);
\draw[arr] (pos) -- node[yes]{正项} (ratio);
\draw[arr] (pos) -- node[no]{交错} (alt);
\draw[arr] (ratio) -- (ratioR);
\draw[arr] (alt) -- (leibniz);
\draw[arr] (ratioR) -- (final);
\draw[arr] (leibniz) -- (final);
\end{tikzpicture}
\end{document}
